\documentclass[10pt]{article}

\usepackage[utf8]{inputenc}

\usepackage[T1]{fontenc}

\usepackage{amsmath}

\usepackage{amsfonts}

\usepackage{amssymb}

\usepackage[version=4]{mhchem}

\usepackage{stmaryrd}

\usepackage{mathrsfs}



\begin{document}

ные представления В. г., где каждому вращению соответствует не одна унитарная матрица, а две различающиеся знаком матрицы. Двузначным (спинорным) представлениям $S O(3)$ соответствуют истинные представления накрывающей группы $S U(2)$.



Произведение неприводимых представлений В. г. не является неприводимым, но может быть разложено в прямую сумму неприводимых представлений. Коэффициенты этого разложения (Клебша - Гордана коэффициенты ) используют в квантовой механике при вычислении матричных элементов различных операторов и при построении волновых функций составных систем.



Вращение на малый угол можно представить в виде



\[

\hat{M}=1+\hat{T}_{12} \varphi_{12}+\hat{T}_{13} \varphi_{13}+\hat{T}_{23} \varphi_{23},

\]



где $\hat{M}$ - матрица вращения в нек-ром представлении, $\varphi_{m n}$ - малые углы поворота в трех независимых плоскостях, a $\hat{T}_{m n}$ - фиксированные матрицы, к-рые в данном представлении называются ге н е р а т о р а м и В. г. В квантовой механике генераторы В. г. имеют наглядный физич. смысл и совпадают с операторами углового момента $\hat{\boldsymbol{J}}$. Некоммутативность В.г. отражается в том факте, что коммутатор $\left[\hat{J}_{m}, \hat{J}_{n}\right]$ отличен от нуля при $m \neq n$.



Отметим, что в нек-рых физич. задачах находят применение и группы $O(n)$ с $n>3$. Так, группа $O(4)$ оказалась полезной при классификации состояний атома водорода, в теории гравитации интерес представляет связанная с группой $O(5)$ de Cummepa групnа, при попытках построения единой квантовой теории поля используют В. г. высших размерностей вплоть до $O(32)$.



Лит.: [1] Л юбарски й Г.Я., Теория групп и ее применение в физике, М., 1958; [2] Гельфанд И. М., Минлос Р.А., Шапи ро 3.Я., Представления группы врашений и группы Лоренца и их применения, М., 1958; [3] Ющ и с А. П., Ле в и н с он И. Б., В а н а га с В. В., Математический аппарат теории момента количества движения, Вильнюс, 1960; [4] П е т р аш е нь М. И., Тр и фо но в Е.Д., Применение теории групп в квантовой механике, М., 1967.



ВРЕМЕНИПОДОБНОЕ ПРОСТРАНСТВО-ВРЕМЯ см. Сингулярности пространства-времени.



BPEMEНИПОДОБНЫЙ BEKTOP - четырехмерный вектор в пространстве-времени специальной теории относительности (Минковского пространстве-времени), квадрат временной компоненты которого больше суммы квадратов пространственных компонент. Квадрат длины В. В. $(A)$ в метрике Минковского положителен:



\[

(A)^{2}=A^{\mu} A_{\mu}=\left(A^{0}\right)^{2}-\left(A^{1}\right)^{2}-\left(A^{2}\right)^{2}-\left(A^{3}\right)^{2} .

\]



Здесь $A^{0}-$ временная, $A^{1}, A^{2}, A^{3}-$ пространственные компоненты $(\mu=0,1,2,3)$. С помошью Лоренца преобразований последние могут быть обращены в нуль, то есть существует система отсчета, в к-рой данный В. в. характеризуется единственной, временной компонентой $A^{0^{\prime}}$. Поскольку величина $(A)^{2}$ инвариантна при преобразованиях Лоренца, $A^{0^{\prime}}=\sqrt{(A)^{2}}$. Очевидно, любой В. в. остается времениподобным при переходе в произвольно движушуюся инерциальную систему отсчета.



Важным примером В. в. в релятивистской механике является вектор четырехмерной скорости частицы с ненулевой массой покоя:



\[

u^{\mu}=d x^{\mu} / d s

\]



\begin{itemize}

  \item касательный вектор к мировой линии $x^{\mu}(s)$ частицы $(s-$ интервал). Этот вектор в метрике Минковского имеет единичную длину, $u^{2}=1$, а система отсчета, в к-рой его пространственные компоненты равны нулю, является системой покоя частицы. В этой системе направление $u^{\mu}$ совпадает с направлением оси времени. С В. в. $u^{\mu}$ связан соотношением пропорциональности другой В. в.- четырехмерный импульс $p^{\mu}=m u^{\mu}(m-$ масса частицы; используется система единиц, в к-рой скорость света $c=1$ ). Временная компонента этого вектора равна полной энергии \& частицы с учетом энергии покоя, $p^{2}=p^{\mu} p_{\mu}=\xi^{2}-p^{2}=m^{2}(\boldsymbol{p}-$ трехмерный импульс).

\end{itemize}



В пространстве-времени Минковского времениподобным будет любой вектор, лежащий внутри светового конуса, вершина к-рого совмещена с его началом. Такой В. в. соединяет точки, отвечающие событиям, к-рые могут быть причинно связаны между собой. Соответствуюший интервал (длина этого вектора) также называется врем енипо д о б н Ы м. BPEMEНHОЕ СРЕДНЕЕ - см. Усреднения метод, Эргодическая теорема.



BРЕМЕННОЕ УРАВНЕНИЕ ШРЕДИНГЕРА - см. Шредингера уравнение.



BPEМЕННОЙ РЯД - см. Случайная функция, Случайный процесс.



BPEMЯ ХИЗНИ - характеристика резонанса (неспектральной особенности) в квантовой теории рассеяния, определяюшая скорость распада резонансного состояния. Если матрица рассеяния $S(k)$, определенная для пары гамильтонианов $H, H_{0}$, действующих в гильбертовом пространстве $\mathscr{\varkappa}$, имеет полюс $k_{0}$ на нефизич. листе $\left(\operatorname{Im} k_{0}<0, k_{0}^{2}=E\right)$ римановой поверхности энергий $E$, то в р е е нем ж и зн и резонанса $k_{0}$ называется величина



\[

\tau_{0}=\left|\operatorname{Im} k_{0}\right|^{-1} \text {. }

\]



Эта величина показывает, с какой скоростью происходит размывание волнового пакета в окрестности рассеивателя. Если $p$ - проектор на состояния, сосредоточенные в окрестности рассеивателя, то норма функции резонансного состояния $U_{0}^{\text {рез }}$ в гильбертовом пространстве $\mathscr{N}$ экспоненциально убывает с течением времени $t$ :



\[

\left\|P U_{0}^{\text {pe3 }}\right\|(t)=\left\|P U_{0}^{\text {pe } 3}\right\|_{(t=0)} \exp \left(-\tau_{0}^{-1} t\right) .

\]



Обратная величина $\Gamma_{0}=\tau_{0}^{-1}$ называется ши и и ной ре зон а н с а.



Лum.: [1] Р и д М., С а й м о н Б., Методы современной математической физики, т. 4. Анализ операторов, пер. с англ., М., 1982.



Ю. Б. Мельников



BTOPAЯ KВАДРАТИЧНАЯ ФOPMA пове рхности, радиус-вектор текущей точки которой задан параметрически $\boldsymbol{r}=\boldsymbol{r}\left(u^{1}, u^{2}\right),-$ форма вида



\[

d \boldsymbol{r} d \boldsymbol{m}=b_{i j} d u^{i} d u^{j}

\]



где $\boldsymbol{m}-$ единичная нормаль к поверхности. В качестве нормали обычно берут



то есть



\[

\boldsymbol{m}=\left[\frac{\partial \boldsymbol{r}}{\partial u^{1}}, \frac{\partial \boldsymbol{r}}{\partial u^{2}}\right] /\left|\left[\frac{\partial \boldsymbol{r}}{\partial u^{1}}, \frac{\partial \boldsymbol{r}}{\partial u^{2}}\right]\right|,

\]



\[

b_{i j}=\left(\frac{\partial \boldsymbol{r}}{\partial u^{1}}, \frac{\partial \boldsymbol{r}}{\partial u^{2}}, \frac{\partial^{2} \boldsymbol{r}}{\partial u^{i} \partial u^{j}}\right) \gamma^{-1 / 2},

\]



где $\gamma=\operatorname{det}\left\|g_{i j}\right\|, g_{i j}-$ первая квадратичная форма.



Часто используют классич. обозначения:



\[

b_{11}=L, b_{12}=M, b_{22}=N .

\]



В. В. Трофимов.



BTOPAЯ KPAEBAЯ ЗАДАЧА - см. Краевая задача, Неймана задача.



BТОРИЧНОЕ KВАНТОВАНИЕ - метоД рассмотрения квантовой системы, при котором роль независимых переменных играет число частиц в заданном состоянии. В.к. возникло при рассмотрении нерелятивистских систем, состоящих из тождественных частиц. Для бозе-частиц (подчиняюшихся статистике Бозе - Эйнштейна) метод В. к. развит в 1927 П. Дираком (P. Dirac) и в том же году - П. Иорданом (P. Jordan) и О. Клейном (O. Klein), для фермичастиц (подчиняюшихся статистике Ферми - Дирака) -





\end{document}